% Abstract text, shared by main.tex (ICGA) and arxiv.tex.
The strongest poker agents approximate an equilibrium and do not adapt, so they leave value
unclaimed against imperfect players. We study how much a lightweight agent can gain by
adapting to opponents using only public information: the betting it observes and the cards
revealed at showdown. Our agent extends a Monte Carlo equity player with an opponent model
that shrinks each observed frequency toward a typical player, and with bounded adjustments
to ranges, bluffing, bet sizing and calling. With no evidence about an opponent, it plays
exactly as the base agent. In heads-up No-Limit Texas Hold'em at 100 big blinds, over
200,000 to 300,000 hands per match, the adaptive agent wins more than the base agent against
all six opponents tested: for example 65.5 against 41.4 big blinds per 100 hands (\bbh)
against an opponent who over-folds, and 234 against 105 against one that switches style to
defeat adaptation. It also beats the base agent directly by 8.4~\bbh{} (95\% interval
$\pm$3.4). All opponents are synthetic and the agent is not an equilibrium player; we state
these limits and outline the path to stronger baselines.
